%! Author = lili
%! Date = 7/30/26


\section{Waterman-Eggert Example}\label{sec:waterman-eggert-example}

\subsubsection*{Input}
\begin{figure}[H]
    \centering
    \begin{tabular}{c || c | c | c | c | c |}
        i $\backslash$ j & \phantom{G} & G & C & T &  A   \\ \hline \hline
        & 0&0 &0 &0 &0  \\ \hline
        G&0 & 3 & 2 & 1 & 1 \\ \hline
        T&0 & 2 & 4 & 5 & 4 \\ \hline
        A&0 & 1 & 3 & 5 & 8 \\ \hline
        A&0 &1 &2 &4 & 8 \\ \hline
    \end{tabular}
    \label{fig:result}
\end{figure}


\subsubsection*{Take the calculated matrix from Smith-Waterman and set all cells that contribute to the optimal alignmnent to 0.}
\begin{figure}[H]
    \centering
    \begin{tabular}{c || c | c | c | c | c |}
        i $\backslash$ j & \phantom{G} & G & C & T &  A   \\ \hline \hline
        & 0&0 &0 &0 &0  \\ \hline
        G&0 & 0 & 2 & 1 & 1 \\ \hline
        T&0 & 2 & 0 & 5 & 4 \\ \hline
        A&0 & 1 & 3 & 0 & 8 \\ \hline
        A&0 &1 &2 &4 & 0\\ \hline
    \end{tabular}
    \label{fig:result}
\end{figure}

\subsubsection*{Main part - calculation}
$i_start \gets $ 1
\begin{algorithmic}
    \For{$i \gets i_start$ to $n$}
        \For{$j \gets 1$ to $m$}
            \If{(i, j) belongs to optimal aligmnent}
                \State $S(i,j) = 0$
            \Else
                \State $S(i,j) \gets max \begin{cases}
                                             0   \\
                                             S(i - 1, j - 1) + s(x[i], y[j])  \\
                                             S(i - 1, j) - s(x[i], --)  \\
                                             S(i, j - 1) - s(--, y[j])
                \end{cases}$
            \EndIf
        \EndFor
    \EndFor
\end{algorithmic}
\paragraph{i = 1, j = 1:}
\begin{multicols}{2}
(1,1) belongs to optimal aligmnent, also $S(i - 1, j - 1) + s(x[i], y[j]) = 0$ bzw. dieser Weg darf nicht gegangen werden.

Und über den Rest bekommt man auch maximal 0, da:
\[
    S(1,1) \gets max \begin{cases}
                         0   \\
                         S(0, 1) - 1 = 0-1 = -1  \\
                         S(1, 0) - 1 = 0-1 = -1
    \end{cases}
\]
\end{multicols}

\paragraph{i = 1, j = 2:}
\begin{multicols}{2}
    \[
        S(1,2) =
        \max \begin{cases}
                 S(0,1)& + \phantom{-} 1 =0+1 = 1 \\
                 S(0,2)& + -1 = 0-1 = -1 \\
                 S(1,1)& + -1 = 0-1 = -1\\
                 0
        \end{cases} \quad  = 1
    \]
    \begin{figure}[H]
        \centering
        \begin{tabular}{c || c | c | c | c | c |}
            i $\backslash$ j & \phantom{G} & G & C & T &  A   \\ \hline \hline
            & 0&0 &0 &0 &0  \\ \hline
            G&0 & 0 & \newn{1} & 1 & 1 \\ \hline
            T&0 & 2 & 0 & 5 & 4 \\ \hline
            A&0 & 1 & 3 & 0 & 8 \\ \hline
            A&0 &1 &2 &4 & 0\\ \hline
        \end{tabular}
        \label{fig:result}
    \end{figure}
\end{multicols}


\paragraph{i = 1, j = 3:}
\begin{multicols}{2}
    \[
        S(1,3) =
        \max \begin{cases}
                 S(0,2)& + \phantom{-}1 = 0+1 = 1 \\
                 S(0,3)& + -1 = 0-1 = -1 \\
                 S(1,2)& + -1 = 1-1 = 0\\
                 0
        \end{cases} \quad  = 1
    \]
    \begin{figure}[H]
        \centering
        \begin{tabular}{c || c | c | c | c | c |}
            i $\backslash$ j & \phantom{G} & G & C & T &  A   \\ \hline \hline
            & 0&0 &0 &0 &0  \\ \hline
            G&0 & 0 & 1 & 1 & 1 \\ \hline
            T&0 & 2 & 0 & 5 & 4 \\ \hline
            A&0 & 1 & 3 & 0 & 8 \\ \hline
            A&0 &1 &2 &4 & 0\\ \hline
        \end{tabular}
        \label{fig:result}
    \end{figure}
\end{multicols}

\paragraph{i = 1, j = 4:}
\begin{multicols}{2}
    \[
        S(1,4) =
        \max \begin{cases}
                 S(0,3)& + \phantom{-} 1 =0+1 = 1 \\
                 S(0,4)& + -1 = 0-1 = -1 \\
                 S(1,3)& + -1 = 1-1 = 0\\
                 0
        \end{cases} \quad  = 1
    \]
    \begin{figure}[H]
        \centering
        \begin{tabular}{c || c | c | c | c | c |}
            i $\backslash$ j & \phantom{G} & G & C & T &  A   \\ \hline \hline
            & 0&0 &0 &0 &0  \\ \hline
            G&0 & 0 & 1 & 1 & 1 \\ \hline
            T&0 & 2 & 0 & 5 & 4 \\ \hline
            A&0 & 1 & 3 & 0 & 8 \\ \hline
            A&0 &1 &2 &4 & 0\\ \hline
        \end{tabular}
        \label{fig:result}
    \end{figure}
\end{multicols}


\paragraph{i = 2, j = 1:}
\begin{multicols}{2}
    \[
        S(2,1) =
        \max \begin{cases}
                 S(1,0)& + \phantom{-}1 = 0+1 = 1 \\
                 S(1,1)& + -1 = 0-1 = -1 \\
                 S(2,0)& + -1 = 0-1 = -1\\
                 0
        \end{cases} \quad  = 1
    \]
    \begin{figure}[H]
        \centering
        \begin{tabular}{c || c | c | c | c | c |}
            i $\backslash$ j & \phantom{G} & G & C & T &  A   \\ \hline \hline
            & 0&0 &0 &0 &0  \\ \hline
            G&0 & 0 & 1 & 1 & 1 \\ \hline
            T&0 & \newn{1} & 0 & 5 & 4 \\ \hline
            A&0 & 1 & 3 & 0 & 8 \\ \hline
            A&0 &1 &2 &4 & 0\\ \hline
        \end{tabular}
        \label{fig:result}
    \end{figure}
\end{multicols}

\paragraph{i = 2, j = 2:}
\begin{multicols}{2}
(2,2) belongs to optimal aligmnent , also $S(i - 1, j - 1) + s(x[i], y[j]) = 0$ bzw. dieser Weg darf nicht gegangen werden.

Und über den Rest bekommt man auch maximal 0, da:
\[
    S(2,2) \gets max \begin{cases}
                         0   \\
                         S(1, 2) - 1 = 0-1 = -1  \\
                         S(2, 1) - 1 = 0-1 = -1
    \end{cases}
\]
\end{multicols}

\paragraph{i = 2, j = 3:}
\begin{multicols}{2}
    \[
        S(2,3) =
        \max \begin{cases}
                 S(1,2)& + \phantom{-} 3 =1+3 = 4 \\
                 S(1,3)& + -1 = 1-1 = 0 \\
                 S(2,2)& + -1 = 0-1 = -1\\
                 0
        \end{cases} \quad  = 4
    \]
    \begin{figure}[H]
        \centering
        \begin{tabular}{c || c | c | c | c | c |}
            i $\backslash$ j & \phantom{G} & G & C & T &  A   \\ \hline \hline
            & 0&0 &0 &0 &0  \\ \hline
            G&0 & 0 & 1 & 1 & 1 \\ \hline
            T&0 & 1 & 0 & \newn{4} & 4 \\ \hline
            A&0 & 1 & 3 & 0 & 8 \\ \hline
            A&0 &1 &2 &4 & 0\\ \hline
        \end{tabular}
        \label{fig:result}
    \end{figure}
\end{multicols}

\paragraph{i = 2, j = 4:}
\begin{multicols}{2}
    \[
        S(2,4) =
        \max \begin{cases}
                 S(1,3)& + \phantom{-}1 = 1+1 = 2 \\
                 S(1,4)& + -1 = 1-1 = 0 \\
                 S(2,3)& + -1 = 4-1 = 3\\
                 0
        \end{cases} \quad  = 3
    \]
    \begin{figure}[H]
        \centering
        \begin{tabular}{c || c | c | c | c | c |}
            i $\backslash$ j & \phantom{G} & G & C & T &  A   \\ \hline \hline
            & 0&0 &0 &0 &0  \\ \hline
            G&0 & 0 & 1 & 1 & 1 \\ \hline
            T&0 & 1 & 0 & 4 & \newn{3} \\ \hline
            A&0 & 1 & 3 & 0 & 8 \\ \hline
            A&0 &1 &2 &4 & 0\\ \hline
        \end{tabular}
        \label{fig:result}
    \end{figure}
\end{multicols}


\paragraph{i = 3, j = 1:}
\begin{multicols}{2}
    \[
        S(3,1) =
        \max \begin{cases}
                 S(2,0)& + \phantom{-} 1 =0+1 = 1 \\
                 S(2,1)& + -1 = 1-1 = 0 \\
                 S(3,0)& + -1 = 0-1 = -1\\
                 0
        \end{cases} \quad  = 1
    \]
    \begin{figure}[H]
        \centering
        \begin{tabular}{c || c | c | c | c | c |}
            i $\backslash$ j & \phantom{G} & G & C & T &  A   \\ \hline \hline
            & 0&0 &0 &0 &0  \\ \hline
            G&0 & 0 & 1 & 1 & 1 \\ \hline
            T&0 & 1 & 0 & 4 & 3 \\ \hline
            A&0 & 1 & 3 & 0 & 8 \\ \hline
            A&0 &1 &2 &4 & 0\\ \hline
        \end{tabular}
        \label{fig:result}
    \end{figure}
\end{multicols}

\paragraph{i = 3, j = 2:}
\begin{multicols}{2}
    \[
        S(3,2) =
        \max \begin{cases}
                 S(2,1)& + \phantom{-} 1 = 1+1 = 2 \\
                 S(2,2)& + -1 = 1-1 = -1 \\
                 S(3,1)& + -1 = 1-1 = 0\\
                 0
        \end{cases} \quad  = 2
    \]
    \begin{figure}[H]
        \centering
        \begin{tabular}{c || c | c | c | c | c |}
            i $\backslash$ j & \phantom{G} & G & C & T &  A   \\ \hline \hline
            & 0&0 &0 &0 &0  \\ \hline
            G&0 & 0 & 1 & 1 & 1 \\ \hline
            T&0 & 1 & 0 & 4 & 3 \\ \hline
            A&0 & 1 & \newn{2} & 0 & 8 \\ \hline
            A&0 &1 &2 &4 & 0\\ \hline
        \end{tabular}
        \label{fig:result}
    \end{figure}
\end{multicols}


\paragraph{i = 3, j = 3:}
\begin{multicols}{2}
(3,3) belongs to optimal aligmnent $\Rightarrow S(3,3) = 0$

Aber über den Rest bekommt man 3, da:
\[
    S(3,3) \gets max \begin{cases}
                         0   \\
                         S(2, 3) - 1 = 4-1 = 3  \\
                         S(3, 2) - 1 = 2-1 = 1
    \end{cases}
\]

\begin{figure}[H]
    \centering
    \begin{tabular}{c || c | c | c | c | c |}
        i $\backslash$ j & \phantom{G} & G & C & T &  A   \\ \hline \hline
        & 0&0 &0 &0 &0  \\ \hline
        G&0 & 0 & 1 & 1 & 1 \\ \hline
        T&0 & 1 & 0 & 4 & 3 \\ \hline
        A&0 & 1 & 2 & \newn{3} & 8\\ \hline
        A&0 &1 &2 &4 & 0\\ \hline
    \end{tabular}
    \label{fig:result}
\end{figure}
\end{multicols}


\paragraph{i = 3, j = 4:}
\begin{multicols}{2}
    \[
        S(3,4) =
        \max \begin{cases}
                 S(2,3)& + \phantom{-} 3 = 4+3= 7 \\
                 S(2,4)& + -1 = 3-1 = 2 \\
                 S(3,3)& + -1 = 3-1 = 2\\
                 0
        \end{cases} \quad  = 7
    \]
    \begin{figure}[H]
        \centering
        \begin{tabular}{c || c | c | c | c | c |}
            i $\backslash$ j & \phantom{G} & G & C & T &  A   \\ \hline \hline
            & 0&0 &0 &0 &0  \\ \hline
            G&0 & 0 & 1 & 1 & 1 \\ \hline
            T&0 & 1 & 0 & 4 & 3 \\ \hline
            A&0 & 1 & 2 & 3 & \newn{7}\\ \hline
            A&0 &1 &2 &4 & 0\\ \hline
        \end{tabular}
        \label{fig:result}
    \end{figure}
\end{multicols}
\paragraph{i = 4, j = 1:}
\begin{multicols}{2}
    \[
        S(4,1) =
        \max \begin{cases}
                 S(3,0)& + \phantom{-} 1 = 0+1 = 1 \\
                 S(3,1)& + -1 = 1-1 = 0 \\
                 S(4,0)& + -1 = 0-1 = -1\\
                 0
        \end{cases} \quad  = 1
    \]
    \begin{figure}[H]
        \centering
        \begin{tabular}{c || c | c | c | c | c |}
            i $\backslash$ j & \phantom{G} & G & C & T &  A   \\ \hline \hline
            & 0&0 &0 &0 &0  \\ \hline
            G&0 & 0 & 1 & 1 & 1 \\ \hline
            T&0 & 1 & 0 & 4 & 3 \\ \hline
            A&0 & 1 & 2 & 3 & 7\\ \hline
            A&0 &1 &2 &4 & 0\\ \hline
        \end{tabular}
        \label{fig:result}
    \end{figure}
\end{multicols}

\paragraph{i = 4, j = 2:}
\begin{multicols}{2}
    \[
        S(4,2) =
        \max \begin{cases}
                 S(3,1)& + \phantom{-}1 = 1+1 = 2 \\
                 S(3,2)& + -1 = 2-1 = 1 \\
                 S(4,1)& + -1 = 1-1 = 0\\
                 0
        \end{cases} \quad  = 2
    \]
    \begin{figure}[H]
        \centering
        \begin{tabular}{c || c | c | c | c | c |}
            i $\backslash$ j & \phantom{G} & G & C & T &  A   \\ \hline \hline
            & 0&0 &0 &0 &0  \\ \hline
            G&0 & 0 & 1 & 1 & 1 \\ \hline
            T&0 & 1 & 0 & 4 & 3 \\ \hline
            A&0 & 1 & 2 & 3 & 7\\ \hline
            A&0 &1 &2 &4 & 0\\ \hline
        \end{tabular}
        \label{fig:result}
    \end{figure}
\end{multicols}

\paragraph{i = 4, j = 3:}
\begin{multicols}{2}
    \[
        S(4,3) =
        \max \begin{cases}
                 S(3,2)& + \phantom{-} 2+1 = 3 \\
                 S(3,3)& + -1 = 3-1 = 2 \\
                 S(4,2)& + -1 = 2-1 = 1\\
                 0
        \end{cases} \quad  = 3
    \]
    \begin{figure}[H]
        \centering
        \begin{tabular}{c || c | c | c | c | c |}
            i $\backslash$ j & \phantom{G} & G & C & T &  A   \\ \hline \hline
            & 0&0 &0 &0 &0  \\ \hline
            G&0 & 0 & 1 & 1 & 1 \\ \hline
            T&0 & 1 & 0 & 4 & 3 \\ \hline
            A&0 & 1 & 2 & 3 & 7\\ \hline
            A&0 &1 &2 & \newn{3} & 0\\ \hline
        \end{tabular}
        \label{fig:result}
    \end{figure}
\end{multicols}

\paragraph{i = 4, j = 4:}
\begin{multicols}{2}
(4,4) belongs to optimal aligmnent

Aber über den Rest bekommt man 6, da:
\[
    S(4,4) \gets max \begin{cases}
                         0   \\
                         S(3, 4) - 1 = 7-1 = 6  \\
                         S(4, 3) - 1 = 3-1 = 2
    \end{cases}
\]

\begin{figure}[H]
    \centering
    \begin{tabular}{c || c | c | c | c | c |}
        i $\backslash$ j & \phantom{G} & G & C & T &  A   \\ \hline \hline
        & 0&0 &0 &0 &0  \\ \hline
        G&0 & 0 & 1 & 1 & 1 \\ \hline
        T&0 & 1 & 0 & 4 & 3 \\ \hline
        A&0 & 1 & 2 & 3 & 7\\ \hline
        A&0 &1 &2 & 3 & \newn{6} \\ \hline
    \end{tabular}
    \label{fig:result}
\end{figure}

\end{multicols}

\subsubsection*{Result}

\begin{figure}[H]
    \centering
    \begin{tabular}{c || c | c | c | c | c |}
        i $\backslash$ j & \phantom{G} & G & C & T &  A   \\ \hline \hline
        & 0&0 &0 &0 &0  \\ \hline
        G&0 & 0 & 1 & 1 & 1 \\ \hline
        T&0 & 1 & 0 & 4 & 3 \\ \hline
        A&0 & 1 & 2 & 3 & 7\\ \hline
        A&0 &1 &2 & 3 & 6\\ \hline
    \end{tabular}
    \label{fig:result}
\end{figure}



\subsection*{Links}
\begin{itemize}
    \item \url{https://www.compbio.dundee.ac.uk/papers/water/node5.html}
    \item \url{https://w.wiki/ZAL}
    \item \url{https://rna.informatik.uni-freiburg.de/Teaching/index.jsp?toolName=Smith-Waterman}
\end{itemize}

WARNUNG: man kann zwar bei eggert das erste Aligment nicht mehr nehmen, aber es könnte dort trtzdem noch ein Indel sein.

% Siehe Foto, da müsste ganz unten rechts eine 4 sein.